% GENERATED FILE. Edit canonical lesson frontmatter, then run npm run generate:books.
% canonical-entries: 53

\chapter*{Glossary}
\addcontentsline{toc}{chapter}{Glossary}
\markboth{Glossary}{Glossary}

This glossary collects every word and phrase taught in the book. Pronunciation appears when it differs from the written form; chapter numbers show where each entry is introduced.

\noindent\begin{minipage}[t]{\linewidth}
\raggedright
\textbf{\gu{આંખ}}\enspace\emph{āṅkh}\par
\small eye — a straight cousin of the English word, and a nasal mark that is not the neuter tell you might expect\par
\footnotesize Introduced in Chapter 16.
\end{minipage}\par\medskip

\noindent\begin{minipage}[t]{\linewidth}
\raggedright
\textbf{\gu{આવવું}}\enspace\emph{āvvũ}\par
\small to come — from a root meaning “reach,” and a cousin of the word copula\par
\footnotesize Introduced in Chapter 11.
\end{minipage}\par\medskip

\noindent\begin{minipage}[t]{\linewidth}
\raggedright
\textbf{\gu{બહેન}}\enspace\emph{bahen}\par
\small sister — which, unlike its brother, is not the same word as the English one at all\par
\footnotesize Introduced in Chapter 15.
\end{minipage}\par\medskip

\noindent\begin{minipage}[t]{\linewidth}
\raggedright
\textbf{\gu{ભાઈ}}\enspace\emph{bhāī}\par
\small brother — the plainest possible cousin of the English word for the very same person\par
\footnotesize Introduced in Chapter 15.
\end{minipage}\par\medskip

\noindent\begin{minipage}[t]{\linewidth}
\raggedright
\textbf{\gu{ચા}}\enspace\emph{chā}\par
\small tea — a word that took the overland road out of China, so it rhymes with Hindi and Persian rather than with English\par
\footnotesize Introduced in Chapter 14.
\end{minipage}\par\medskip

\noindent\begin{minipage}[t]{\linewidth}
\raggedright
\textbf{\gu{દૂધ}}\enspace\emph{dūdh}\par
\small milk — inherited straight from Sanskrit, and a root English keeps in an unexpected word\par
\footnotesize Introduced in Chapter 14.
\end{minipage}\par\medskip

\noindent\begin{minipage}[t]{\linewidth}
\raggedright
\textbf{\gu{એક} \gu{બે} \gu{ત્રણ} \gu{ચાર} \gu{પાંચ}}\enspace\emph{ek be traṇ chār pā̃ch}\par
\small one to five — the two numbers that don't look like their neighbours'\par
\footnotesize Introduced in Chapter 10.
\end{minipage}\par\medskip

\noindent\begin{minipage}[t]{\linewidth}
\raggedright
\textbf{\gu{ગમવું}}\enspace\emph{gamvũ}\par
\small to like — the chapter's payoff, a verb that will not let you be its subject, born as the passive of \textquotedblleft{}to go\textquotedblright{}\par
\footnotesize Introduced in Chapter 13.
\end{minipage}\par\medskip

\noindent\begin{minipage}[t]{\linewidth}
\raggedright
\textbf{\gu{હોવું}}\enspace\emph{hovũ}\par
\small to be — and the neuter ending that names every Gujarati verb\par
\footnotesize Introduced in Chapter 11.
\end{minipage}\par\medskip

\noindent\begin{minipage}[t]{\linewidth}
\raggedright
\textbf{\gu{જાણવું}}\enspace\emph{jāṇvũ}\par
\small to know — the word English has been carrying all along\par
\footnotesize Introduced in Chapter 11.
\end{minipage}\par\medskip

\noindent\begin{minipage}[t]{\linewidth}
\raggedright
\textbf{\gu{જવું}}\enspace\emph{javũ}\par
\small to go — and the Sanskrit y- that Prakrit turned into j-\par
\footnotesize Introduced in Chapter 11.
\end{minipage}\par\medskip

\noindent\begin{minipage}[t]{\linewidth}
\raggedright
\textbf{\gu{જોવું}}\enspace\emph{jovũ}\par
\small to see — one vowel sign away from “to go,” and an origin still argued over\par
\footnotesize Introduced in Chapter 11.
\end{minipage}\par\medskip

\noindent\begin{minipage}[t]{\linewidth}
\raggedright
\textbf{\gu{કાન}}\enspace\emph{kān}\par
\small ear — a word Sanskrit itself could not settle the origin of, and one that changed gender on the way to Gujarati\par
\footnotesize Introduced in Chapter 16.
\end{minipage}\par\medskip

\noindent\begin{minipage}[t]{\linewidth}
\raggedright
\textbf{\gu{ખાવું}}\enspace\emph{khāvũ}\par
\small to eat — the verb whose family reaches its Indian sisters and stops there\par
\footnotesize Introduced in Chapter 11.
\end{minipage}\par\medskip

\noindent\begin{minipage}[t]{\linewidth}
\raggedright
\textbf{\gu{કુટુંબ}}\enspace\emph{kuṭumb}\par
\small family — the whole group, and a word that arrived in Gujarati already finished rather than growing there\par
\footnotesize Introduced in Chapter 15.
\end{minipage}\par\medskip

\noindent\begin{minipage}[t]{\linewidth}
\raggedright
\textbf{\gu{લખવું}}\enspace\emph{lakhvũ}\par
\small to write — the chapter's payoff, on a root that meant to scratch, and a family that stops before English again\par
\footnotesize Introduced in Chapter 12.
\end{minipage}\par\medskip

\noindent\begin{minipage}[t]{\linewidth}
\raggedright
\textbf{\gu{લેવું}}\enspace\emph{levũ}\par
\small to take — the shortest verb yet, and a family tree that is disputed rather than empty\par
\footnotesize Introduced in Chapter 13.
\end{minipage}\par\medskip

\noindent\begin{minipage}[t]{\linewidth}
\raggedright
\textbf{\gu{મદદ} \gu{કરવી}}\enspace\emph{madad karvī}\par
\small to help — an Arabic word that came by way of Persian, and the one place the verb ending stops being neuter\par
\footnotesize Introduced in Chapter 13.
\end{minipage}\par\medskip

\noindent\begin{minipage}[t]{\linewidth}
\raggedright
\textbf{\gu{મિત્ર}}\enspace\emph{mitra}\par
\small friend — a word this chapter will link to the very phrase you used to ask for water\par
\footnotesize Introduced in Chapter 15.
\end{minipage}\par\medskip

\noindent\begin{minipage}[t]{\linewidth}
\raggedright
\textbf{\gu{મોઢું}}\enspace\emph{moḍhũ}\par
\small mouth — the clearest example yet of Gujarati's neuter tell, the nasalized -ũ you have been hearing on every verb\par
\footnotesize Introduced in Chapter 16.
\end{minipage}\par\medskip

\noindent\begin{minipage}[t]{\linewidth}
\raggedright
\textbf{\gu{નાક}}\enspace\emph{nāk}\par
\small nose — the fourth body word, closing the chapter by running the whole face back through gender, sound, and certainty\par
\footnotesize Introduced in Chapter 16.
\end{minipage}\par\medskip

\noindent\begin{minipage}[t]{\linewidth}
\raggedright
\textbf{\gu{નમસ્તે}}\enspace\emph{namaste}\par
\small hello / goodbye (namaste)\par
\footnotesize Introduced in Chapter 1.
\end{minipage}\par\medskip

\noindent\begin{minipage}[t]{\linewidth}
\raggedright
\textbf{\gu{પાણી}}\enspace\emph{pāṇī}\par
\small water — Gujarati's first noun for something you would actually ask for, and the phrase that asks for it politely\par
\footnotesize Introduced in Chapter 14.
\end{minipage}\par\medskip

\noindent\begin{minipage}[t]{\linewidth}
\raggedright
\textbf{\gu{પૂછવું}}\enspace\emph{pūchhvũ}\par
\small to ask — the verb English kept only for praying, and a p that stayed put because of where it stood\par
\footnotesize Introduced in Chapter 13.
\end{minipage}\par\medskip

\noindent\begin{minipage}[t]{\linewidth}
\raggedright
\textbf{\gu{રોટલી}}\enspace\emph{roṭlī}\par
\small flatbread — the fourth request, an honest dead end for its origin, and the letter -ī does not always mean what you'd guess\par
\footnotesize Introduced in Chapter 14.
\end{minipage}\par\medskip

\noindent\begin{minipage}[t]{\linewidth}
\raggedright
\textbf{\gu{સમજવું}}\enspace\emph{samajvũ}\par
\small to understand — built on a root that meant \textquotedblleft{}to wake up,\textquotedblright{} and the cluster it lost on the way\par
\footnotesize Introduced in Chapter 12.
\end{minipage}\par\medskip

\noindent\begin{minipage}[t]{\linewidth}
\raggedright
\textbf{\gu{વાંચવું}}\enspace\emph{vā̃chvũ}\par
\small to read — which once meant \textquotedblleft{}to make something speak,\textquotedblright{} and a nasal that lives in the middle of the word\par
\footnotesize Introduced in Chapter 12.
\end{minipage}\par\medskip

\noindent\begin{minipage}[t]{\linewidth}
\raggedright
\textbf{\gu{વિચારવું}}\enspace\emph{vichārvũ}\par
\small to think — a verb Gujarati built for itself, on a root that meant to wander\par
\footnotesize Introduced in Chapter 12.
\end{minipage}\par\medskip

\noindent\begin{minipage}[t]{\linewidth}
\raggedright
\textbf{\gu{આભાર}}\par
\small thank you (ābhār)\par
\footnotesize Introduced in Chapter 2.
\end{minipage}\par\medskip

\noindent\begin{minipage}[t]{\linewidth}
\raggedright
\textbf{\gu{આવજો}}\par
\small goodbye (lit. \textquotedblleft{}{[}please{]} come {[}again{]}\textquotedblright{})\par
\footnotesize Introduced in Chapter 2.
\end{minipage}\par\medskip

\noindent\begin{minipage}[t]{\linewidth}
\raggedright
\textbf{\gu{કેમ}}\par
\small how (also \textquotedblleft{}why\textquotedblright{})\par
\footnotesize Introduced in Chapter 7.
\end{minipage}\par\medskip

\noindent\begin{minipage}[t]{\linewidth}
\raggedright
\textbf{\gu{કામ} \gu{કરવું}}\par
\small to work (lit. \textquotedblleft{}work-do\textquotedblright{})\par
\footnotesize Introduced in Chapter 9.
\end{minipage}\par\medskip

\noindent\begin{minipage}[t]{\linewidth}
\raggedright
\textbf{\gu{કાલે} \gu{મળીશું}}\par
\small see you tomorrow (and the \gu{કાલ} puzzle)\par
\footnotesize Introduced in Chapter 8.
\end{minipage}\par\medskip

\noindent\begin{minipage}[t]{\linewidth}
\raggedright
\textbf{\gu{છે}}\par
\small is (the copula)\par
\footnotesize Introduced in Chapter 6.
\end{minipage}\par\medskip

\noindent\begin{minipage}[t]{\linewidth}
\raggedright
\textbf{\gu{તું} / \gu{તમે}}\par
\small you (familiar / respectful)\par
\footnotesize Introduced in Chapter 6.
\end{minipage}\par\medskip

\noindent\begin{minipage}[t]{\linewidth}
\raggedright
\textbf{\gu{તમે} \gu{કેમ} \gu{છો}?}\par
\small how are you? (respectful)\par
\footnotesize Introduced in Chapter 7.
\end{minipage}\par\medskip

\noindent\begin{minipage}[t]{\linewidth}
\raggedright
\textbf{\gu{તમને} \gu{મળીને} \gu{આનંદ} \gu{થયો}}\par
\small pleased to meet you (lit. \textquotedblleft{}meeting you, joy happened\textquotedblright{})\par
\footnotesize Introduced in Chapter 6.
\end{minipage}\par\medskip

\noindent\begin{minipage}[t]{\linewidth}
\raggedright
\textbf{\gu{તમારું} \gu{નામ} \gu{શું} \gu{છે}?}\par
\small what's your name?\par
\footnotesize Introduced in Chapter 6.
\end{minipage}\par\medskip

\noindent\begin{minipage}[t]{\linewidth}
\raggedright
\textbf{\gu{નામ}}\par
\small name\par
\footnotesize Introduced in Chapter 6.
\end{minipage}\par\medskip

\noindent\begin{minipage}[t]{\linewidth}
\raggedright
\textbf{\gu{પાછા}}\par
\small again, back\par
\footnotesize Introduced in Chapter 8.
\end{minipage}\par\medskip

\noindent\begin{minipage}[t]{\linewidth}
\raggedright
\textbf{\gu{પાછા} \gu{મળીશું}}\par
\small we'll meet again (goodbye)\par
\footnotesize Introduced in Chapter 8.
\end{minipage}\par\medskip

\noindent\begin{minipage}[t]{\linewidth}
\raggedright
\textbf{\gu{બોલવું}}\par
\small to speak\par
\footnotesize Introduced in Chapter 9.
\end{minipage}\par\medskip

\noindent\begin{minipage}[t]{\linewidth}
\raggedright
\textbf{\gu{મજા}}\par
\small enjoyment, fun — and the reply \textquotedblleft{}\gu{હું} \gu{મજામાં} \gu{છું}\textquotedblright{}\par
\footnotesize Introduced in Chapter 7.
\end{minipage}\par\medskip

\noindent\begin{minipage}[t]{\linewidth}
\raggedright
\textbf{\gu{મારું}}\par
\small my (neuter; also \gu{મારો} m. / \gu{મારી} f.)\par
\footnotesize Introduced in Chapter 6.
\end{minipage}\par\medskip

\noindent\begin{minipage}[t]{\linewidth}
\raggedright
\textbf{\gu{મારું} \gu{નામ} … \gu{છે}}\par
\small my name is…\par
\footnotesize Introduced in Chapter 6.
\end{minipage}\par\medskip

\noindent\begin{minipage}[t]{\linewidth}
\raggedright
\textbf{\gu{મળીશું}}\par
\small (we) will meet\par
\footnotesize Introduced in Chapter 8.
\end{minipage}\par\medskip

\noindent\begin{minipage}[t]{\linewidth}
\raggedright
\textbf{\gu{રહેવું}}\par
\small to live, to stay\par
\footnotesize Introduced in Chapter 9.
\end{minipage}\par\medskip

\noindent\begin{minipage}[t]{\linewidth}
\raggedright
\textbf{\gu{વાંધો} \gu{નહીં}}\par
\small no problem / it's okay / you're welcome\par
\footnotesize Introduced in Chapter 7.
\end{minipage}\par\medskip

\noindent\begin{minipage}[t]{\linewidth}
\raggedright
\textbf{\gu{શું}}\par
\small what\par
\footnotesize Introduced in Chapter 6.
\end{minipage}\par\medskip

\noindent\begin{minipage}[t]{\linewidth}
\raggedright
\textbf{\gu{સારું}}\par
\small good, okay, fine (sārũ)\par
\footnotesize Introduced in Chapter 2.
\end{minipage}\par\medskip

\noindent\begin{minipage}[t]{\linewidth}
\raggedright
\textbf{\gu{હું}}\par
\small I\par
\footnotesize Introduced in Chapter 7.
\end{minipage}\par\medskip

\noindent\begin{minipage}[t]{\linewidth}
\raggedright
\textbf{\gu{હા} / \gu{ના}}\par
\small yes / no (hā / nā)\par
\footnotesize Introduced in Chapter 2.
\end{minipage}\par\medskip

\noindent\begin{minipage}[t]{\linewidth}
\raggedright
\textbf{\gu{હું} \gu{ગુજરાતી} \gu{બોલું} \gu{છું}}\par
\small I speak Gujarati\par
\footnotesize Introduced in Chapter 9.
\end{minipage}\par\medskip
